\section[Ejemplo]{Ejemplo: categoría con clivaje definida por un funtor
$\cal{E}^\circ\to\Cat$; categorías escindidas sobre~$\cal{E}$}
\label{VI.9}
Supongamos dado un funtor
$$
\ifthenelse{\boolean{orig}}
{\phi\colon\cal{E}^\circ\to\Cat,}
{\varphi\colon\cal{E}^\circ\to\Cat,}
$$
Este define entonces un pseudofuntor al poner
$$
\ifthenelse{\boolean{orig}}
{\cal{F}(S)=\phi(S),}
{\cal{F}(S)=\varphi(S),}
\quad f^*=\varphi(f),\quad c_{f,g}=\id_{(fg)^*}
$$
\marginpar{180}%
Por tanto, la construcción del número precedente nos da una categoría
$\cal{F}$ con clivaje sobre~$\cal{E}$, llamada asociada al funtor
$\varphi$. Para que una categoría con clivaje sobre~$\cal{E}$ sea
isomorfa a una categoría con clivaje definida por un funtor
$\varphi:\cal{E}^\circ\to\Cat$, es necesario y suficiente,
manifiestamente, que satisfaga las condiciones
$$
(fg)^*=g^*f^*,\quad c_{f,g}=\id_{(fg)^*}\quoi.
$$
En términos del conjunto $K$ de los morfismos de transporte, esto significa
sencillamente que \emph{el compuesto de dos morfismos de transporte es un
morfismo de transporte}. Un clivaje de una categoría $\cal{F}$
sobre~$\cal{E}$ que satisface la condición precedente se llama una
\emph{escisión}
\index{escision@escisión|hyperpage}%
de~$\cal{F}$ sobre~$\cal{E}$, y una categoría $\cal{F}$ sobre~$\cal{E}$
dotada de una escisión se llama una \emph{categoría escindida
sobre~$\cal{E}$}.
\index{escindida (categoria)@escindida (categoría)|hyperpage}%
Se trata, pues, de un caso particular de la noción de categoría con
clivaje. La categoría de las categorías escindidas sobre~$\cal{E}$ es,
por tanto, equivalente a~$\SheafHom(\cal{E}^\circ,\Cat)$. Obsérvese que
una categoría escindida sobre~$\cal{E}$ es, con mayor razón, una
categoría con clivaje sobre~$\cal{E}$.

Si $\cal{F}$ es una categoría fibrada sobre~$\cal{E}$, no siempre existe
una escisión de~$\cal{F}$. Supongamos, por ejemplo, que
$\Ob\cal{E}$ y $\Ob\cal{F}$ se reducen cada uno a un elemento y que los
endomorfismos de estos elementos forman grupos $E$ y $F$,
respectivamente, de modo que el funtor de proyección $p$ viene dado por un
homomorfismo de grupos $p\colon F\to E$, sobreyectivo porque $p$ es un
fibrado. Se comprueba entonces inmediatamente que el conjunto de los
clivajes de~$\cal{F}$ sobre~$\cal{E}$ está en correspondencia biunívoca
con el conjunto de las aplicaciones $s\colon E\to F$ tales que
$ps=\id_E$ (es decir, el conjunto de los «sistemas de representantes» de
las clases módulo el subgrupo $G$, núcleo del homomorfismo sobreyectivo
$p\colon F\to E$). Un clivaje es una escisión si y solo si $s$ es un
homomorfismo de grupos. Decir que existe una escisión significa, pues, que
la extensión de grupos $F$ de~$E$ por~$G$ es escindida, lo cual, cuando
$G$ es conmutativo, se expresa por la anulación de cierta clase de
cohomología en~$\H^2(E,G)$ (donde se considera $G$ como un grupo sobre el
que actúa~$E$).

Supongamos, sin embargo, que $\cal{F}$ es una categoría fibrada
sobre~$\cal{E}$ tal que las
\marginpar{181}
$\cal{F}_S$ son categorías \emph{rígidas}, es decir, el grupo de
automorfismos de todo objeto de~$\cal{F}_S$ se reduce a la identidad. Es
entonces fácil probar que $\cal{F}$ admite una escisión sobre~$\cal{E}$.
En efecto, se observa primero que la cuestión de la existencia de una
escisión no cambia si se sustituye $\cal{F}$ por una categoría
$\cal{E}$-equivalente, lo que nos reduce en este caso a aquel en que las
$\cal{F}_S$ son categorías rígidas \emph{y reducidas} (es decir, dos
objetos isomorfos de~$\cal{F}_S$ son idénticos). Pero si $G$ es una
categoría rígida y reducida, todo isomorfismo entre dos funtores $H\to G$
(donde $H$ es una categoría cualquiera) es una identidad. Se sigue que si
$\cal{F}$ es una categoría fibrada sobre~$\cal{E}$ cuyas categorías fibra
son rígidas y reducidas, entonces existe un clivaje \emph{único} de
$\cal{F}$ sobre~$\cal{E}$, que es necesariamente una escisión. Por tanto,
$\cal{F}$ es isomorfa a la categoría definida por un funtor
$\varphi\colon \cal{E}^\circ\to \Cat$ tal que las $\varphi(S)$ sean
categorías rígidas y
\ifthenelse{\boolean{orig}}
{discretas}
{reducidas}
, y el funtor $\varphi$ queda determinado salvo isomorfismo.
